This section is split into locked local diagnostics and pending Qwen results.

\subsection{Benign Curves Overstate HLC Durability}

On benign relearn pools, HLC-SG appeared strongest in the local MVP.
It had the lowest immediate resurrection and highest benign survival AUC.
This result motivated the stricter conditional protocol.

\begin{table}[h]
\centering
\caption{Benign local MVP main-track result. HLC appears strongest before immediate matching and answer-bearing stress.}
\label{tab:benign-main}
\begin{tabular}{lccc}
\toprule
Method & Immediate $r@0$ & AUC survival & $h50$ \\
\midrule
NPO & 0.833 & 0.167 & 0 observed \\
Strong NPO & 0.222 & 0.778 & $>32$ censored \\
NPO + paraphrases & 0.167 & 0.833 & $>32$ censored \\
Syntactic diversification & 0.444 & 0.556 & mixed \\
HLC-SG & 0.056 & 0.944 & $>32$ censored \\
\bottomrule
\end{tabular}
\end{table}


\subsection{Answer-Bearing Stress Changes the Conclusion}

Under immediate-matched answer-bearing stress, current HLC variants did not beat the strongest local baselines.
The held-out stress table is the clearest locked diagnostic result.

\begin{table}[h]
\centering
\caption{Held-out answer-bearing stress result from the locked local MVP. HLC beats NPO + paraphrases but not immediate-matched Strong NPO or syntactic diversification.}
\label{tab:heldout-stress}
\begin{tabular}{lccccc}
\toprule
Method & $r@0$ & $c_{32}$ & $c_{128}$ & $c_{512}$ & Cond. AUC$_{512}$ \\
\midrule
HLC-SG & 0.056 & 0.037 & 0.685 & 0.978 & 0.351 \\
NPO + paraphrases & 0.167 & 0.733 & 1.000 & 1.000 & 0.057 \\
Strong NPO matched & 0.056 & 0.000 & 0.263 & 0.978 & 0.410 \\
Syntactic diversification matched & 0.056 & 0.056 & 0.404 & 0.981 & 0.480 \\
\bottomrule
\end{tabular}
\end{table}


\subsection{GPT-2 LoRA Smoke Confirms the Larger-Model Path}

The GPT-2 LoRA lane validates the training and evaluation path but does not establish a method win.
K2 improves some short-horizon signals over K1, but still trails Strong NPO by $t=128$.

\begin{table}[h]
\centering
\caption{GPT-2 LoRA long-horizon diagnostic. K2 improves over K1 at $t=128$ but does not beat Strong NPO.}
\label{tab:gpt2-long}
\begin{tabular}{lccccc}
\toprule
Method & $r@0$ & $c_{32}$ & AUC$_{32}$ & $c_{128}$ & AUC$_{128}$ \\
\midrule
HLC-SG K1 GPT-2 & 0.000 & 0.278 & 0.920 & 0.833 & 0.529 \\
HLC-SG K2 GPT-2 & 0.000 & 0.278 & 0.920 & 0.611 & 0.584 \\
Strong NPO & 0.000 & 0.333 & 0.906 & 0.556 & 0.588 \\
\bottomrule
\end{tabular}
\end{table}


\subsection{Qwen3-0.6B Seeded Long-Horizon Gate}

The Qwen3-0.6B lane is the active decision gate for a tuned HLC K4 candidate against immediate-matched constrained gradient ascent.
The unseeded/provisional long-512 result should not be used for paper claims.
The seeded rerun must populate this table.

\begin{table}[h]
\centering
\caption{Qwen3-0.6B seeded long-horizon decision gate. Fill from the seeded Qwen report after the rerun completes.}
\label{tab:qwen-pending}
\small
\begin{tabular}{lccccc}
\toprule
Method & $r@0$ & $c_{32}$ & $c_{128}$ & $c_{512}$ & Cond. AUC$_{512}$ \\
\midrule
HLC K4 & \todo{fill} & \todo{fill} & \todo{fill} & \todo{fill} & \todo{fill} \\
Matched GA & \todo{fill} & \todo{fill} & \todo{fill} & \todo{fill} & \todo{fill} \\
\bottomrule
\end{tabular}
\end{table}


\paragraph{Current interpretation.}
If tuned HLC beats matched gradient ascent on conditional resurrection or conditional AUC after the seeded rerun, the paper can include a positive Qwen diagnostic result.
If it loses, the paper should remain a rigorous negative result centered on the evaluation protocol.
